%=============================================================
% Paper 7: Mueller (2017), AER
%=============================================================
\chapter{Mueller (2017)}

\begin{paperbox}{Separations, Sorting, and Cyclical Unemployment}
\textbf{Author:} Andreas I.\ Mueller (Columbia GSB)\\[3pt]
\textbf{Journal:} \textit{American Economic Review} 107(7): 2081--2107\\[3pt]
\textbf{Year:} 2017 \quad\textbf{Field:} Labour economics; Macroeconomics\\[3pt]
\textbf{Classification:} Empirical + structural (calibration)
\end{paperbox}

%------------------------------------------------------------
\section{Research Question}

Does the composition of the pool of unemployed workers shift systematically
over the business cycle, and in which direction---toward higher- or lower-wage
workers? What drives this compositional shift---separations or job-finding
rates? Can standard or modified theories of unemployment account for the
observed pattern?

%------------------------------------------------------------
\section{Motivation}

It has long been known that the pool of \emph{employed} workers shifts toward
higher-wage workers in recessions (Solon, Barsky, Parker 1994), because
low-wage workers are more likely to lose jobs. One might expect the
\emph{unemployed} pool to shift toward lower-wage workers. Mueller documents
that this intuition is wrong: the unemployed pool shifts toward
\emph{high}-wage workers in recessions. The finding has profound implications
for hiring incentives (compositional changes affect the average productivity
of job applicants), for policy evaluation (ALMP participants are higher-quality
in recessions), and for the unemployment volatility puzzle (Shimer 2005).

%------------------------------------------------------------
\section{Main Contribution}

\begin{enumerate}
  \item \textbf{A new empirical fact:} The pre-displacement wage of unemployed
    workers rises strongly in recessions. A 1~pp increase in the unemployment
    rate is associated with a 1.5~pp increase in the average wage \emph{rank}
    of the unemployed ($\rho=0.67$, CPS~ORG; 0.73, March~CPS). This holds
    for raw wages, wage ranks, and Mincer residuals, and appears in every
    recession since 1962.
  \item \textbf{Mechanism:} The shift is driven almost entirely by
    \emph{separations}---high-wage workers have substantially more cyclical
    separation rates; job-finding rate cyclicality is similar across wage
    groups.
  \item \textbf{Theory evaluation:} Standard DMP models fail. Two extensions
    succeed: (i)~lower variance of match productivity for high-ability workers;
    (ii)~credit cash-flow constraints in recessions.
\end{enumerate}

%------------------------------------------------------------
\section{Relationship to the Literature}

Extends Solon--Barsky--Parker (1994) to the unemployed pool, showing both
pools can shift in the \emph{same} direction (high-wage) in recessions. Challenges
Pries (2008) who assumes the unemployed pool shifts toward low-ability workers.
Adds a new challenge to the unemployment volatility puzzle (Shimer 2005): the
high-ability unemployed raise vacancy returns, \emph{dampening} the aggregate
hiring response to adverse shocks.

%------------------------------------------------------------
\section{Theoretical Framework and Economic Mechanism}

\subsection*{Accounting Framework}

The fraction of group~$i$ in the unemployed pool satisfies:
\begin{equation}
  \phi^U_{it} = \phi^L_{it}\,\frac{U_{it}}{U_t},
  \label{eq:mueller_phi}
\end{equation}
where $\phi^L_{it}$ is group $i$'s labour force share and $U_{it}$ is group
$i$'s unemployment rate. Using the Elsby--Michaels--Solon (2009) flow
decomposition, compositional changes in the unemployed pool decompose
(approximately) as:
\begin{equation}
  d\phi^U_{it} \approx \phi^U_{it}(1-\phi^U_{it})\,\bigl[
    (1-U^{ss}_{it})(d\ln s_{it}-d\ln f_{it})
    - (1-U^{ss}_{jt})(d\ln s_{jt}-d\ln f_{jt})
  \bigr],
  \label{eq:mueller_decomp}
\end{equation}
where $s_{it}$ and $f_{it}$ are group-$i$ separation and job-finding rates.
This shows compositional changes are driven by the \emph{relative} log changes
in flow rates, weighted by steady-state employment rates.

\subsection*{Why Standard Models Fail}

In the baseline DMP model, low-wage workers are more likely to be near the
separation threshold, so an adverse aggregate shock pushes more low-wage
workers into unemployment---opposite to the data.

\subsection*{Extension 1: Lower Match Variance for High-Ability Workers}

If the variance of match-specific productivity shocks is lower for high-ability
workers, the density of matches at the separation threshold is
\emph{higher} for high-ability workers (their matches are more tightly
clustered near the threshold). Aggregate shocks that shift the threshold
cause a larger mass of high-ability matches to cross it, generating higher
cyclicality of high-wage separations. Critically, this calibration also
preserves near-equal cyclicality of job-finding rates across groups.

\subsection*{Extension 2: Credit Cash-Flow Constraints}

Firms facing cash-flow constraints in recessions may separate from workers
whose current match productivity is negative even if the NPV of the match is
positive. Marginal matches with high-ability workers generate more negative
current cash flows (workers paid above current productivity in anticipation
of future returns), making them more sensitive to credit tightening.

%------------------------------------------------------------
\section{Data}

\begin{itemize}
  \item \textbf{CPS Outgoing Rotation Groups (ORG):} Private-sector workers
    aged 16--64, 1980--2012; $N=1{,}203{,}543$; 79,463 experienced $\geq$1
    month of unemployment in interview months 5--8.
  \item \textbf{CPS March Supplement:} 1962--2012 (backward-looking wages;
    full 50-year coverage; no attrition issues).
  \item \textbf{Monthly CPS:} 1978--2012 (demographic composition of
    unemployed; no wage data).
  \item \textbf{NLSY79:} 1979--2012; $N=6{,}923$; 193,467 yearly labour
    force status observations; used to distinguish permanent vs.\ transitory
    wage components.
\end{itemize}

%------------------------------------------------------------
\section{Empirical Strategy}

\begin{enumerate}
  \item Compute average log previous wage, wage rank (ordinal), and Mincer
    residual for the unemployed pool in each year. HP-filter ($\lambda=100$)
    and correlate with the HP-filtered aggregate unemployment rate.
  \item Decompose compositional shifts into components by predicted wage
    (experience, education, gender, marital status, race, state, industry,
    occupation) using OLS regressions (Table~1).
  \item Estimate group-specific $(s_{it}, f_{it})$ from matched CPS panels
    and apply the decomposition in equation~(\ref{eq:mueller_decomp}) to
    isolate the separation vs.\ job-finding margin.
  \item Calibrate baseline and extended DMP models to match average separation
    rates by wage group; compare predicted compositional shifts to data.
\end{enumerate}

%------------------------------------------------------------
\section{Identification Strategy}

\textbf{The main empirical claim is correlational, not causal.} The paper
documents a robust pattern across multiple data sources and specifications.
The Mincer residual specification (controlling for demographics, industry,
occupation, year) attempts to isolate the residual component. The NLSY79
analysis provides evidence that compositional changes are associated with
\emph{permanent} worker fixed effects rather than transitory match quality
effects.

Key identification challenge: the composition of workers who lose their jobs
could reflect true changes in separation behaviour \emph{or} changes in the
composition of at-risk workers. The paper addresses this through the wage rank
specification (immune to wage compression) and the Mincer residual (controlling
for observables), but cannot fully rule out unobservable worker selection into
high-wage jobs.

%------------------------------------------------------------
\section{Main Results}

\subsection*{New Empirical Fact (Figures 3--4, Table~1)}

\begin{center}
\begin{tabular}{lcc}
\toprule
\textbf{Wage measure} & \textbf{CPS ORG} & \textbf{March CPS} \\
\midrule
$\rho$: avg log previous wage $\times$ unemp.\ rate & 0.60 & 0.59 \\
$\rho$: avg wage \emph{rank} $\times$ unemp.\ rate & 0.67 & 0.73 \\
$\rho$: avg Mincer residual $\times$ unemp.\ rate & 0.44 & 0.60 \\
\midrule
1~pp $\uparrow$ unemp.\ rate $\Rightarrow$ $\Delta$ avg log wage & $+2.77\%$ & $+2.59\%$ \\
1~pp $\uparrow$ unemp.\ rate $\Rightarrow$ $\Delta$ avg wage rank & $+1.5$~pp & $+1.1$~pp \\
1~pp $\uparrow$ unemp.\ rate $\Rightarrow$ $\Delta$ avg Mincer residual & $+0.75\%$ & $+0.95\%$ \\
\bottomrule
\end{tabular}
\end{center}

The pattern appears in \emph{every recession since 1962} and is robust to all
specifications.

\subsection*{Observable Drivers (Table~1)}

All observable components shift in the same direction: the unemployed in
recessions are more experienced, more educated, more likely to be male,
married, white, and from high-wage industries and occupations. Industry
accounts for $\approx$25\% of total compositional shift.

\subsection*{Separation vs.\ Job-Finding}

The compositional shift is \emph{almost entirely driven by separations}.
High-wage workers have substantially more cyclical log separation rates;
the cyclicality of log job-finding rates is similar across wage groups.

\subsection*{Model Results}

\begin{itemize}
  \item \textbf{Standard DMP:} Generates shifts toward \emph{low}-ability
    workers---wrong direction.
  \item \textbf{Extension~1 (lower match variance):} Generates shifts toward
    high-ability workers; correct direction for all moments.
  \item \textbf{Extension~2 (credit constraints):} Correct direction for
    recessions with financial stress; internally inconsistent for recessions
    without credit disruption.
\end{itemize}

%------------------------------------------------------------
\section{Economic Magnitude and Interpretation}

A 1~pp rise in unemployment shifts the average unemployed worker from (say)
the 50th to the 51.5th percentile of the wage distribution. Over a typical
recession ($\approx$3~pp unemployment rise), the average unemployed worker
is at about the 54.5th percentile. The estimated dampening of aggregate
hiring responses due to this compositional shift is up to $3.6\times$ the
steady-state elasticity (Online Appendix~G.1)---a quantitatively important
challenge for the unemployment volatility puzzle.

%------------------------------------------------------------
\section{Mechanisms}

The preferred explanation is \textbf{lower variance of match-specific
productivity for high-ability workers}:
\begin{itemize}
  \item High-ability workers have matches clustered more tightly near the
    separation threshold.
  \item Aggregate shocks that shift the threshold by a given amount cause
    a larger mass of high-ability matches to cross it.
  \item This generates higher cyclicality of separations without asymmetric
    cyclicality of job-finding rates.
\end{itemize}

The \textbf{credit constraint} mechanism is plausible but has internal
consistency problems: it predicts no compositional shift in recessions without
clear financial market disruption (e.g.\ 1970--71), but the data show the
shift in every recession back to 1962.

%------------------------------------------------------------
\section{Robustness Checks}

\begin{itemize}
  \item Three CPS data sources (ORG, March, monthly) and NLSY79 give
    consistent results.
  \item Robust to: wage rank (immune to wage compression); Mincer residual
    (controls for observables); exclusion of long-term unemployed.
  \item NLSY79 distinguishes permanent worker effects from transitory match
    effects: shift is associated with permanent quality.
  \item Pattern holds in every recession 1962--2012.
\end{itemize}

%------------------------------------------------------------
\section{Limitations}

\begin{enumerate}
  \item \textbf{Long-term unemployment exclusion:} CPS ORG excludes $>$12
    months unemployed; becomes substantial post-2008 (up to 31.2\%).
  \item \textbf{Indirect mechanism test:} Cannot directly test whether
    match productivity variance differs across worker types---a maintained
    calibration assumption.
  \item \textbf{Credit constraint test indirect:} Chodorow-Reich (2014)-based
    test is promising but not definitive.
\end{enumerate}

%------------------------------------------------------------
\section{Policy Implications}

\begin{itemize}
  \item \textbf{ALMP evaluation:} Higher-ability unemployed in recessions
    $\Rightarrow$ programmes more effective in recessions (Card, Kluve,
    Weber 2015).
  \item \textbf{UI design:} Cyclical shift toward high-wage unemployed
    means UI benefits (fixed fraction of prior wage) are \emph{larger in
    absolute terms} during recessions---an automatic stabiliser amplification.
  \item \textbf{Hiring incentives:} High-ability unemployed raise returns
    to vacancy posting $\Rightarrow$ partial automatic stabiliser on the
    hiring side (dampens vacancy collapse).
\end{itemize}

%------------------------------------------------------------
\section{Overall Assessment}

Mueller (2017) is a clean, careful, and important empirical paper establishing
a new stylised fact with exemplary robustness: three data sources, 50 years
of data, multiple specifications, and a clear accounting framework. The theory
section is more provisional but usefully narrows the space of viable
explanations. The most important contribution is the implication for the
unemployment volatility puzzle: compositional shifts toward high-ability
workers in recessions \emph{dampen} aggregate hiring responses.

\textbf{Quality: Very good. Important new fact, careful empirics, sound
theoretical interpretation.}

%------------------------------------------------------------
\section{Relevance to My Research}

Relevant for any RDC project studying worker heterogeneity in unemployment
dynamics. For Canadian work: LEED or matched LFS-style data could be used to
test whether the Canadian unemployed pool also shifts toward high-wage workers
in recessions, and whether this effect is amplified in resource provinces
during commodity busts (where separated workers are predominantly from
high-wage extraction jobs).

%------------------------------------------------------------
\section{Potential Research Extensions}

\begin{itemize}
  \item Test the Mueller fact for Canada using matched LFS data; does the
    Canadian unemployed pool shift toward high-wage workers, and is this
    amplified in resource provinces during commodity busts?
  \item COVID-19 (primarily destroying low-wage service jobs) likely produced
    the \emph{opposite} compositional shift---a natural out-of-sample test
    of the mechanism's asymmetry.
\end{itemize}

%------------------------------------------------------------
\section{Research Gaps}

\begin{itemize}
  \item Cannot distinguish the credit constraint from the lower match-variance
    mechanism; a firm-level test using credit supply shocks (e.g.\ bank
    failures) would be needed.
  \item COVID-19 provides an obvious out-of-sample test of whether the
    Mueller mechanism operates symmetrically---the composition of unemployment
    should shift toward \emph{low}-wage workers.
\end{itemize}

%------------------------------------------------------------
\section{Methodological Lessons}

\begin{enumerate}
  \item Equation~(\ref{eq:mueller_phi}) is essential: changes in the share of
    a group among the unemployed depend on \emph{relative} unemployment rate
    changes (not absolute levels), so the higher-volatility unemployment of
    low-wage workers does \emph{not} imply the unemployed pool shifts toward
    low-wage workers.
  \item Wage \emph{ranks} (ordinal) are more informative than levels because
    they are immune to wage compression as an alternative explanation.
  \item The Elsby--Michaels--Solon flow decomposition (equation
    \ref{eq:mueller_decomp}) cleanly attributes compositional shifts to
    separations vs.\ job-finding, using log rates (not levels) to account
    for different average flow rates across groups.
\end{enumerate}

%------------------------------------------------------------
\section*{Five-Sentence Summary}

Mueller (2017) documents a new fact using 50 years of CPS data: in recessions,
the pool of unemployed shifts toward workers with \emph{higher} pre-displacement
wages, with a 1~percentage-point rise in unemployment associated with a 1.5~pp
increase in the average wage rank of the unemployed---a pattern that holds in
every recession since 1962, for raw wages, wage ranks, and Mincer residuals.
The compositional shift is almost entirely driven by \emph{separations}: high-
wage workers have substantially more cyclical separation rates, while job-
finding rate cyclicality is similar across wage groups. Standard DMP models
generate shifts in the \emph{opposite} direction; two extensions succeed---
lower variance of match productivity for high-ability workers and credit cash-
flow constraints---with the former being the more broadly applicable
explanation. The finding implies that the unemployed pool is higher quality in
recessions, raising returns to vacancy posting and potentially \emph{dampening}
the aggregate hiring response to adverse shocks---an additional challenge for
the unemployment volatility puzzle. Compositional changes in the unemployed
pool are substantially larger in magnitude than the corresponding changes in
the employed pool, implying standard composition-bias corrections are
insufficient for statistics related to the unemployed.

\vspace{6pt}
\begin{quickref}
\begin{tabularx}{\linewidth}{@{}lX@{}}
\textbf{Key contribution} &
  New empirical fact (robust across 3 data sources, 50 years): unemployed
  pool shifts toward high-wage workers in recessions, driven almost entirely
  by the higher cyclicality of separations for high-wage workers. \\[4pt]
\textbf{Key weakness} &
  The two successful theoretical extensions cannot be empirically
  distinguished; the credit constraint mechanism is internally inconsistent
  for recessions without financial stress. \\[4pt]
\textbf{Most interesting gap} &
  COVID-19 likely produced the \emph{opposite} compositional shift (low-wage
  service workers primarily displaced)---a natural out-of-sample test of
  whether the Mueller mechanism is asymmetric. \\[4pt]
\textbf{Publishable extension} &
  Replicate Mueller's decomposition for Canada using matched LFS data; test
  whether the shift toward high-wage unemployed is amplified in resource
  provinces during commodity busts, and assess whether it dampens hiring
  responses in those regions.
\end{tabularx}
\end{quickref}
